\documentclass[11pt,a4paper]{book}

\usepackage[margin=2.5cm]{geometry}
\usepackage{amsmath,amssymb,amsthm,bm,mathtools}
\usepackage{tikz}
\usetikzlibrary{calc}
\usepackage{listings}
\usepackage{booktabs}
\usepackage[hidelinks]{hyperref}

\newtheorem{proposition}{Proposition}[chapter]
\theoremstyle{definition}
\newtheorem{exercise}{Exercise}[chapter]
\theoremstyle{remark}
\newtheorem*{remark}{Remark}

\newcommand{\dd}{\mathrm{d}}
\newcommand{\pb}[2]{\{#1,#2\}}
\DeclareMathOperator{\sgn}{sgn}

\lstset{
  language=C,
  basicstyle=\ttfamily\small,
  keywordstyle=\bfseries,
  commentstyle=\itshape,
  frame=single,
  numbers=left,
  numberstyle=\tiny,
  breaklines=true,
  columns=fullflexible
}

\renewcommand{\bibname}{References}

\begin{document}

\setcounter{chapter}{0}
\chapter{The Classical Double Pendulum}
\label{ch:double-pendulum}

\begin{quote}\itshape
A system with two degrees of freedom, two masses, two rods and one uniform gravitational field:
the smallest mechanical system in which every major formalism of classical mechanics has something
to say, and the smallest in which they all fail to give a closed-form solution.
\end{quote}

\section{Introduction and Scope}

The simple pendulum is the first nontrivial problem of classical mechanics: one degree of freedom,
one conserved quantity, and a solution in terms of elliptic functions. Attach a second pendulum to the
bob of the first and the character of the problem changes completely. The system now has two degrees of
freedom but only one \emph{known} independent integral of motion (the energy), so the Liouville--Arnold
theorem no longer applies. The motion is regular at low energy and chaotic at high energy, and the
transition between the two regimes is a textbook illustration of the Kolmogorov--Arnold--Moser theory.

The double pendulum is also an ideal vehicle for comparing the different formulations of classical
mechanics. In this chapter we derive its equations of motion in the following ways:

\begin{enumerate}
  \item Newton's second law with constraint forces (Section~\ref{sec:newton});
  \item d'Alembert's principle and Lagrange multipliers in Cartesian coordinates (Section~\ref{sec:cartesian});
  \item the Euler--Lagrange equations in generalized coordinates (Section~\ref{sec:lagrange});
  \item Hamilton's canonical equations, Poisson brackets and the symplectic structure (Section~\ref{sec:hamilton});
  \item the Hamilton--Jacobi equation, separability, and action--angle variables in the small-oscillation limit
        (Section~\ref{sec:hj});
  \item further formulations: Hamilton's variational principle, the Jacobi--Maupertuis principle, the
        Gibbs--Appell equations (Section~\ref{sec:others}).
\end{enumerate}
We then analyse small oscillations (Section~\ref{sec:small}), the structure of the phase space and the onset of
chaos (Section~\ref{sec:chaos}), and numerical integration (Section~\ref{sec:numerics}). A set of exercises closes
the chapter.

General references for the formalisms used here are Goldstein, Poole and Safko~\cite{goldstein}, Landau and
Lifshitz~\cite{landau} and, for the geometric viewpoint, Arnold~\cite{arnold} and Abraham and
Marsden~\cite{abrahammarsden}. Kane's method is treated in~\cite{kanelevinson}; regular and chaotic Hamiltonian
dynamics in~\cite{lichtenberg}; and geometric numerical integration in~\cite{ghw,sanzserna,leimkuhler}. The
double pendulum itself, as a paradigm of chaos in a simple mechanical system, has been studied numerically and
experimentally in~\cite{shinbrot,levien,stachowiak}.

\section{The Model}
\label{sec:model}

\subsection{Setup and assumptions}

Two point masses $m_1$ and $m_2$ are connected by massless, rigid, inextensible rods of lengths $l_1$ and $l_2$.
The first rod is attached to a fixed frictionless pivot $O$; the second rod joins $m_1$ to $m_2$ through a
frictionless hinge. The motion is confined to a vertical plane in a uniform gravitational field of
magnitude $g$. We write $\mu \equiv m_1+m_2$ for the total mass.

\begin{figure}[ht]
\centering
\begin{tikzpicture}[scale=1.6,>=stealth]
  \coordinate (O) at (0,0);
  \coordinate (A) at (-60:1.5);
  \coordinate (B) at ($(A)+(-35:1.2)$);
  \draw[dashed,gray] (O) -- (0,-2.6);
  \draw[dashed,gray] (A) -- ($(A)+(0,-1.2)$);
  \draw[thick] (O) -- (A) -- (B);
  \fill (O) circle (1.2pt);
  \fill (A) circle (2.5pt) node[right=3pt] {$m_1$};
  \fill (B) circle (3pt) node[right=3pt] {$m_2$};
  \draw (0,-0.7) arc (-90:-60:0.7);
  \node at (0.17,-0.95) {$\theta_1$};
  \draw ($(A)+(0,-0.6)$) arc (-90:-35:0.6);
  \node at ($(A)+(0.28,-0.8)$) {$\theta_2$};
  \node[above left] at ($(O)!0.5!(A)$) {$l_1$};
  \node[above right] at ($(A)!0.5!(B)$) {$l_2$};
  \draw[->] (2.2,-0.3) -- (2.2,-1.1) node[midway,right] {$\bm g$};
  \node[above] at (O) {$O$};
\end{tikzpicture}
\caption{The planar double pendulum. Both angles are measured from the downward vertical, counter-clockwise
positive.}
\label{fig:dp}
\end{figure}

\subsection{Configuration space and degrees of freedom}

Without constraints, the two particles in the plane have four degrees of freedom. The two holonomic constraints
\begin{equation}
  |\bm r_1| = l_1, \qquad |\bm r_2-\bm r_1| = l_2
  \label{eq:constraints}
\end{equation}
leave two. We choose the angles $(\theta_1,\theta_2)$ of Fig.~\ref{fig:dp} as generalized coordinates. The
configuration space is the two-torus $T^2=S^1\times S^1$ and the phase space is its cotangent bundle $T^*T^2$,
of dimension four. Taking the origin at $O$ with the $y$-axis pointing up,
\begin{align}
  x_1 &= l_1\sin\theta_1, & y_1 &= -l_1\cos\theta_1, \label{eq:pos1}\\
  x_2 &= l_1\sin\theta_1 + l_2\sin\theta_2, & y_2 &= -l_1\cos\theta_1 - l_2\cos\theta_2. \label{eq:pos2}
\end{align}
Differentiating,
\begin{align}
  \dot{x}_1 &= l_1\dot\theta_1\cos\theta_1, & \dot{y}_1 &= l_1\dot\theta_1\sin\theta_1,\\
  \dot{x}_2 &= l_1\dot\theta_1\cos\theta_1 + l_2\dot\theta_2\cos\theta_2, &
  \dot{y}_2 &= l_1\dot\theta_1\sin\theta_1 + l_2\dot\theta_2\sin\theta_2,
\end{align}
and therefore, with $\Delta\equiv\theta_1-\theta_2$,
\begin{equation}
  v_1^2 = l_1^2\dot\theta_1^2, \qquad
  v_2^2 = l_1^2\dot\theta_1^2 + l_2^2\dot\theta_2^2 + 2l_1l_2\dot\theta_1\dot\theta_2\cos\Delta .
  \label{eq:v2}
\end{equation}

\section{Newtonian Formulation}
\label{sec:newton}

In the Newtonian approach we treat the rod forces explicitly. Define, for $i=1,2$, the radial and tangential
unit vectors
\begin{equation}
  \bm e_i = (\sin\theta_i,\,-\cos\theta_i), \qquad
  \bm t_i = (\cos\theta_i,\ \sin\theta_i),
\end{equation}
so that $\dot{\bm e}_i=\dot\theta_i\bm t_i$ and $\dot{\bm t}_i=-\dot\theta_i\bm e_i$. Let $T_1$ and $T_2$ be the
tensions in the upper and lower rods, and $\bm g=(0,-g)$. The forces are:
\begin{itemize}
  \item on $m_2$: gravity $m_2\bm g$ and the pull of the lower rod, $-T_2\bm e_2$;
  \item on $m_1$: gravity $m_1\bm g$, the pull of the upper rod toward the pivot, $-T_1\bm e_1$, and the reaction
        of the lower rod, $+T_2\bm e_2$.
\end{itemize}
Newton's second law gives
\begin{align}
  m_1\ddot{\bm r}_1 &= m_1\bm g - T_1\bm e_1 + T_2\bm e_2, \label{eq:newton1}\\
  m_2\ddot{\bm r}_2 &= m_2\bm g - T_2\bm e_2. \label{eq:newton2}
\end{align}
The accelerations follow from \eqref{eq:pos1}--\eqref{eq:pos2}:
\begin{align}
  \ddot{\bm r}_1 &= l_1\big(\ddot\theta_1\bm t_1 - \dot\theta_1^2\bm e_1\big),\\
  \ddot{\bm r}_2 &= \ddot{\bm r}_1 + l_2\big(\ddot\theta_2\bm t_2 - \dot\theta_2^2\bm e_2\big).
\end{align}
We need the scalar products
\begin{equation}
  \bm e_1\!\cdot\!\bm e_2=\bm t_1\!\cdot\!\bm t_2=\cos\Delta,\qquad
  \bm t_1\!\cdot\!\bm e_2=-\sin\Delta,\qquad
  \bm e_1\!\cdot\!\bm t_2=\sin\Delta,
\end{equation}
and $\bm g\!\cdot\!\bm e_i=g\cos\theta_i$, $\bm g\!\cdot\!\bm t_i=-g\sin\theta_i$.

\paragraph{Lower mass.} Projecting \eqref{eq:newton2} on $\bm e_2$ and $\bm t_2$ gives
\begin{align}
  T_2 &= m_2\big[g\cos\theta_2 + l_1\ddot\theta_1\sin\Delta + l_1\dot\theta_1^2\cos\Delta + l_2\dot\theta_2^2\big],
  \label{eq:T2}\\
  l_2\ddot\theta_2 + l_1\ddot\theta_1\cos\Delta - l_1\dot\theta_1^2\sin\Delta + g\sin\theta_2 &= 0.
  \label{eq:eom2}
\end{align}

\paragraph{Upper mass.} Projecting \eqref{eq:newton1} on $\bm t_1$ eliminates $T_1$ (which is orthogonal to the
motion of $m_1$):
\begin{equation}
  m_1 l_1\ddot\theta_1 = -m_1 g\sin\theta_1 - T_2\sin\Delta .
  \label{eq:tangent1}
\end{equation}
Substituting \eqref{eq:T2} for $T_2$ gives the second equation of motion. After a short manipulation (using
\eqref{eq:eom2} to trade $\ddot\theta_2$ for $\ddot\theta_1$) one obtains
\begin{equation}
  \mu\, l_1\ddot\theta_1 + m_2 l_2\ddot\theta_2\cos\Delta + m_2 l_2\dot\theta_2^2\sin\Delta + \mu g\sin\theta_1 = 0 .
  \label{eq:eom1}
\end{equation}
The radial projection of \eqref{eq:newton1} then determines $T_1$ once $\theta_i(t)$ is known. The tensions are
thus a by-product, not an unknown that must be solved for simultaneously, \emph{provided} one projects the
equations onto the directions orthogonal to the constraint forces. This projection is exactly what the Lagrangian
formalism automates.

\section{Cartesian Coordinates and Lagrange Multipliers}
\label{sec:cartesian}

\subsection{d'Alembert's principle}

d'Alembert's principle states that the constraint forces do no virtual work on virtual displacements compatible with the
constraints:
\begin{equation}
  \sum_{k=1}^{2}\big(\bm F_k - m_k\ddot{\bm r}_k\big)\cdot\delta\bm r_k = 0 ,
\end{equation}
where $\bm F_k=m_k\bm g$ are the applied forces. The constraints \eqref{eq:constraints} imply
$\bm r_1\cdot\delta\bm r_1=0$ and $(\bm r_2-\bm r_1)\cdot(\delta\bm r_2-\delta\bm r_1)=0$, so $\delta\bm r_k$ is
not arbitrary. Choosing $\delta\bm r_k=\sum_i(\partial\bm r_k/\partial\theta_i)\,\delta\theta_i$ with independent
$\delta\theta_i$ yields the Lagrange equations of Section~\ref{sec:lagrange}.

\subsection{Multiplier formulation}

Alternatively, keep the redundant Cartesian coordinates $q=(\bm r_1,\bm r_2)\in\mathbb R^4$ and impose the constraints
with multipliers. Write
\begin{equation}
  \phi_1=\tfrac12\big(\bm r_1^2-l_1^2\big)=0,\qquad
  \phi_2=\tfrac12\big((\bm r_2-\bm r_1)^2-l_2^2\big)=0 ,
\end{equation}
and extend the Lagrangian to $L_\lambda=\sum_k\tfrac12 m_k\dot{\bm r}_k^2-\sum_k m_kgy_k-\lambda_1\phi_1-\lambda_2\phi_2$.
Varying with respect to $\bm r_k$ and $\lambda_a$ gives
\begin{align}
  m_1\ddot{\bm r}_1 &= m_1\bm g - \lambda_1\bm r_1 + \lambda_2(\bm r_2-\bm r_1),\\
  m_2\ddot{\bm r}_2 &= m_2\bm g - \lambda_2(\bm r_2-\bm r_1),\\
  \phi_1&=\phi_2=0 .
\end{align}
Comparison with \eqref{eq:newton1}--\eqref{eq:newton2} identifies $T_2=\lambda_2 l_2$ and
$T_1=\lambda_1 l_1 - \lambda_2\,\bm e_1\!\cdot\!(\bm r_2-\bm r_1)$. The multipliers therefore \emph{are} the
(scaled) constraint forces; this is the principal practical advantage of the method.

The price is that the system is now a differential--algebraic equation (DAE). Differentiating $\phi_a$ twice
in time gives algebraic equations for $\lambda_a$ (index-3 DAE reduced to index-1 after two differentiations). In
numerical work one needs constraint stabilization (Baumgarte~\cite{baumgarte}) or projection methods to prevent drift off the
constraint manifold. For the double pendulum, the use of the angles $\theta_i$ avoids the problem altogether.

\section{Lagrangian Formulation}
\label{sec:lagrange}

\subsection{Kinetic and potential energy}

From \eqref{eq:v2}, the kinetic energy $T=\tfrac12m_1v_1^2+\tfrac12m_2v_2^2$ is
\begin{equation}
  T=\tfrac12\mu\,l_1^2\dot\theta_1^2+\tfrac12m_2l_2^2\dot\theta_2^2+m_2l_1l_2\dot\theta_1\dot\theta_2\cos(\theta_1-\theta_2).
  \label{eq:T}
\end{equation}
The potential energy, with zero at the pivot height, is
\begin{equation}
  V=m_1gy_1+m_2gy_2=-\mu g\,l_1\cos\theta_1-m_2g\,l_2\cos\theta_2 .
  \label{eq:V}
\end{equation}
The Lagrangian is $L=T-V$. The kinetic energy is a quadratic form $T=\tfrac12\dot{\bm\theta}^{\!\top}M(\bm\theta)\dot{\bm\theta}$
with the configuration-dependent mass matrix
\begin{equation}
  M(\bm\theta)=
  \begin{pmatrix}
    \mu\,l_1^2 & m_2l_1l_2\cos\Delta\\[2pt]
    m_2l_1l_2\cos\Delta & m_2l_2^2
  \end{pmatrix},
  \qquad
  \det M=m_2l_1^2l_2^2\,\big(m_1+m_2\sin^2\Delta\big).
  \label{eq:M}
\end{equation}
The determinant is strictly positive for $m_1>0$, so $M$ is positive definite everywhere, as it must be for a
physical kinetic energy. (For $m_1\to0$ the determinant vanishes at $\Delta=0,\pi$: the upper bob becomes a
massless joint and the coordinate system degenerates.)

\subsection{Euler--Lagrange equations}

The Euler--Lagrange equations $\frac{\dd}{\dd t}\frac{\partial L}{\partial\dot\theta_i}-\frac{\partial L}{\partial\theta_i}=0$
require the generalized momenta
\begin{align}
  p_1=\frac{\partial L}{\partial\dot\theta_1}&=\mu\,l_1^2\dot\theta_1+m_2l_1l_2\dot\theta_2\cos\Delta,\\
  p_2=\frac{\partial L}{\partial\dot\theta_2}&=m_2l_2^2\dot\theta_2+m_2l_1l_2\dot\theta_1\cos\Delta,
\end{align}
and the generalized forces
\begin{align}
  \frac{\partial L}{\partial\theta_1}&=-m_2l_1l_2\dot\theta_1\dot\theta_2\sin\Delta-\mu g\,l_1\sin\theta_1,\\
  \frac{\partial L}{\partial\theta_2}&=+m_2l_1l_2\dot\theta_1\dot\theta_2\sin\Delta-m_2g\,l_2\sin\theta_2 .
\end{align}
Computing $\dot p_i$ and cancelling the terms proportional to $\dot\theta_1\dot\theta_2\sin\Delta$ yields
\begin{align}
  \mu\,l_1^2\ddot\theta_1+m_2l_1l_2\ddot\theta_2\cos\Delta+m_2l_1l_2\dot\theta_2^2\sin\Delta+\mu g\,l_1\sin\theta_1&=0,
  \label{eq:EL1}\\
  m_2l_2^2\ddot\theta_2+m_2l_1l_2\ddot\theta_1\cos\Delta-m_2l_1l_2\dot\theta_1^2\sin\Delta+m_2g\,l_2\sin\theta_2&=0.
  \label{eq:EL2}
\end{align}
Dividing by $l_1$ and $m_2l_2$ respectively gives \eqref{eq:eom1} and \eqref{eq:eom2}, which were obtained
in Section~\ref{sec:newton} through the Newtonian route. The Lagrangian derivation, however, did not require
knowing the tensions at any stage.

\subsection{Matrix form and explicit accelerations}

In matrix form, the equations of motion read
\begin{equation}
  M(\bm\theta)\,\ddot{\bm\theta}+\Gamma(\bm\theta;\dot{\bm\theta},\dot{\bm\theta})+\nabla V=0,
  \qquad
  \Gamma_i=\sum_{jk}\Gamma_{ijk}\dot\theta_j\dot\theta_k,\quad
  \Gamma_{ijk}=\tfrac12\Big(\frac{\partial M_{ij}}{\partial\theta_k}+\frac{\partial M_{ik}}{\partial\theta_j}-\frac{\partial M_{jk}}{\partial\theta_i}\Big),
  \label{eq:matrixform}
\end{equation}
where $\Gamma_{ijk}$ are the Christoffel symbols of the first kind of the metric $M_{ij}$ on the configuration
torus: the free motion of the system is geodesic motion with respect to the kinetic-energy metric. For the double
pendulum, the only $\theta$-dependence of $M$ is through $\Delta$, so $\Gamma_{\,\cdot\,jk}$ reduces to the two
terms in \eqref{eq:EL1}--\eqref{eq:EL2}, namely $\dot\theta_2^2\sin\Delta$ and $\dot\theta_1^2\sin\Delta$.

Solving the $2\times2$ linear system for the accelerations (by Cramer's rule, using \eqref{eq:M}) gives the
standard explicit form
\begin{align}
  \ddot\theta_1&=\frac{-g(2m_1+m_2)\sin\theta_1-m_2g\sin(\theta_1-2\theta_2)-2m_2\sin\Delta\,\big(l_2\dot\theta_2^2+l_1\dot\theta_1^2\cos\Delta\big)}
  {l_1\big(2m_1+m_2-m_2\cos2\Delta\big)},\label{eq:acc1}\\[4pt]
  \ddot\theta_2&=\frac{2\sin\Delta\,\big(l_1\mu\,\dot\theta_1^2+g\mu\cos\theta_1+l_2m_2\dot\theta_2^2\cos\Delta\big)}
  {l_2\big(2m_1+m_2-m_2\cos2\Delta\big)} .\label{eq:acc2}
\end{align}
Note that $2m_1+m_2-m_2\cos2\Delta=2(m_1+m_2\sin^2\Delta)$, so the denominators are proportional to $\det M$
and never vanish.

\subsection{Symmetries and conserved quantities}

Noether's theorem associates a conserved quantity with each continuous symmetry of $L$. The Lagrangian is
independent of time, so the energy
\begin{equation}
  E=\sum_i\dot\theta_i\frac{\partial L}{\partial\dot\theta_i}-L=T+V
\end{equation}
is conserved. There is \emph{no} other continuous symmetry: the simultaneous rotation $\theta_i\to\theta_i+c$ would
leave $T$ invariant (it depends only on $\Delta$) but it changes $V$ because gravity breaks rotational invariance
about $O$. (In the absence of gravity, $g=0$, the total angular momentum $p_1+p_2$ would be conserved and the system
would be integrable.) With only one integral of motion for two degrees of freedom the system is not Liouville
integrable by this route, and this is the root of its chaotic behaviour.

\begin{remark}
Because $L$ has no ignorable coordinate, Routh's procedure of eliminating cyclic coordinates is of no use here. It is
useful in related problems: the spherical double pendulum, for example, has the azimuthal symmetry about the vertical axis
and the Routhian reduces the number of degrees of freedom from four to two.
\end{remark}

\section{Hamiltonian Formulation}
\label{sec:hamilton}

\subsection{Legendre transform}

The Legendre transformation $H(\bm\theta,\bm p)=\bm p\cdot\dot{\bm\theta}-L$ is well defined because $M$ is
positive definite, which makes $\bm p=M\dot{\bm\theta}$ invertible. Inverting the momenta,
\begin{align}
  \dot\theta_1&=\frac{l_2p_1-l_1p_2\cos\Delta}{l_1^2l_2\,(m_1+m_2\sin^2\Delta)},\label{eq:thd1}\\
  \dot\theta_2&=\frac{\mu\,l_1p_2-m_2l_2p_1\cos\Delta}{m_2l_1l_2^2\,(m_1+m_2\sin^2\Delta)}.\label{eq:thd2}
\end{align}
The Hamiltonian, which is the total energy expressed in canonical variables, is
\begin{equation}
  H=\tfrac12\bm p^{\top}M^{-1}\bm p+V
   =\frac{m_2l_2^2p_1^2+\mu\,l_1^2p_2^2-2m_2l_1l_2\,p_1p_2\cos(\theta_1-\theta_2)}
         {2m_2l_1^2l_2^2\,\big(m_1+m_2\sin^2(\theta_1-\theta_2)\big)}
   -\mu g\,l_1\cos\theta_1-m_2g\,l_2\cos\theta_2 .
  \label{eq:H}
\end{equation}

\subsection{Canonical equations}

Hamilton's equations $\dot\theta_i=\partial H/\partial p_i$, $\dot p_i=-\partial H/\partial\theta_i$ give for the
position equations exactly \eqref{eq:thd1}--\eqref{eq:thd2}. For the momenta, one can differentiate
\eqref{eq:H} directly, but it is much simpler to use the general identity (valid for $H=\tfrac12p^\top M^{-1}p+V$)
$\partial H/\partial\theta_i|_{p}=-\partial L/\partial\theta_i|_{\dot\theta}$. Therefore
\begin{align}
  \dot p_1&=-m_2l_1l_2\,\dot\theta_1\dot\theta_2\sin\Delta-\mu g\,l_1\sin\theta_1,\\
  \dot p_2&=+m_2l_1l_2\,\dot\theta_1\dot\theta_2\sin\Delta-m_2g\,l_2\sin\theta_2,
\end{align}
with $\dot\theta_1,\dot\theta_2$ understood as the functions of $(\bm\theta,\bm p)$ in \eqref{eq:thd1}--\eqref{eq:thd2}.
These four first-order equations form the Hamiltonian system on $\mathbb R^4$ (more precisely on $T^2\times\mathbb R^2$).

\subsection{Poisson brackets and symplectic structure}

With $z=(\theta_1,\theta_2,p_1,p_2)$ and the canonical Poisson bracket
\begin{equation}
  \pb{F}{G}=\sum_{i=1}^{2}\Big(\frac{\partial F}{\partial\theta_i}\frac{\partial G}{\partial p_i}
  -\frac{\partial F}{\partial p_i}\frac{\partial G}{\partial\theta_i}\Big),
\end{equation}
the equations of motion are $\dot z=\pb{z}{H}$ and any observable $F(z,t)$ evolves as
\begin{equation}
  \frac{\dd F}{\dd t}=\pb{F}{H}+\frac{\partial F}{\partial t}.
\end{equation}
The canonical brackets $\pb{\theta_i}{p_j}=\delta_{ij}$ hold, and the flow preserves the symplectic form
$\omega=\sum_i\dd p_i\wedge\dd\theta_i$. Two immediate consequences:

\begin{proposition}[Liouville]
The Hamiltonian flow preserves the phase-space volume $\dd\theta_1\,\dd\theta_2\,\dd p_1\,\dd p_2$.
\end{proposition}
\begin{proof}
The divergence of the vector field is $\sum_i\big(\partial\dot\theta_i/\partial\theta_i+\partial\dot p_i/\partial p_i\big)
=\sum_i\big(\partial^2H/\partial\theta_i\partial p_i-\partial^2H/\partial p_i\partial\theta_i\big)=0$.
\end{proof}

\begin{proposition}[Poisson's theorem and non-integrability]
If $F$ and $G$ are integrals of motion, so is $\pb{F}{G}$. For the double pendulum, the only globally known integral
is $H$.
\end{proposition}
The existence of a \emph{second}, functionally independent, analytic integral $F$ with $\pb{F}{H}=0$ would make the
system Liouville integrable. Such an integral does not exist for generic parameters: the double pendulum has been shown
to be non-integrable in the analytic sense by methods based on the variational equations along particular
solutions and differential Galois theory (the Ziglin and Morales-Ruiz--Ramis approach~\cite{ziglin,moralesruiz,mrr}). The numerical evidence is
the Poincar\'e sections of Section~\ref{sec:chaos}.

\subsection{Canonical transformations}

A transformation $(\bm\theta,\bm p)\to(\bm Q,\bm P)$ is canonical if it preserves $\omega$ or, equivalently, the
fundamental brackets. Generating functions of type 2, $F_2(\bm\theta,\bm P,t)$, give
\begin{equation}
  p_i=\frac{\partial F_2}{\partial\theta_i},\qquad Q_i=\frac{\partial F_2}{\partial P_i},\qquad
  K=H+\frac{\partial F_2}{\partial t}.
  \label{eq:F2}
\end{equation}
The goal of the Hamilton--Jacobi theory is to choose $F_2$ so that $K$ depends on $\bm P$ only; then all $Q_i$
advance linearly in time.

\section{Hamilton--Jacobi Theory}
\label{sec:hj}

\subsection{The Hamilton--Jacobi equation}

Choose $F_2=S(\bm\theta,\bm\alpha,t)$ with constants $\bm\alpha=(\alpha_1,\alpha_2)$ such that the new Hamiltonian
vanishes, $K=0$. Then $\bm P=\bm\alpha$, $\bm Q=\bm\beta$ are constants of motion, and $S$ must satisfy
\begin{equation}
  \frac{\partial S}{\partial t}+H\Big(\bm\theta,\frac{\partial S}{\partial\bm\theta}\Big)=0 .
  \label{eq:HJ}
\end{equation}
For the double pendulum, with $S_i\equiv\partial S/\partial\theta_i$, this is the nonlinear first-order PDE
\begin{equation}
  \frac{\partial S}{\partial t}
  +\frac{m_2l_2^2S_1^2+\mu\,l_1^2S_2^2-2m_2l_1l_2\,S_1S_2\cos\Delta}{2m_2l_1^2l_2^2\,(m_1+m_2\sin^2\Delta)}
  -\mu g\,l_1\cos\theta_1-m_2g\,l_2\cos\theta_2=0 .
  \label{eq:HJexplicit}
\end{equation}
Since $H$ does not depend on $t$, we separate time: $S=W(\bm\theta,\bm\alpha)-Et$ with $\alpha_1=E$, giving the
reduced (time-independent) Hamilton--Jacobi equation
\begin{equation}
  H\Big(\bm\theta,\frac{\partial W}{\partial\bm\theta}\Big)=E .
  \label{eq:HJtimeindep}
\end{equation}
A \emph{complete integral} $W(\theta_1,\theta_2;E,\alpha_2)$, depending on two independent constants, would solve the
motion by quadratures: $\beta_i=\partial W/\partial\alpha_i$ (with $\beta_1=t+t_0$), and $p_i=\partial W/\partial\theta_i$.

\subsection{Separability: why it fails}

Equation \eqref{eq:HJtimeindep} is usually solved by separation of variables, $W=W_1(\theta_1)+W_2(\theta_2)$. Such a
separation works only if the Hamiltonian has a Liouville or St\"ackel form in the chosen coordinates, i.e.\ if
$H$ can be written as
\[
  H=\frac{\sum_i\big(p_i^2/2+U_i(\theta_i)\big)\,(\text{cofactors})}{\det\Phi}.
\]
In \eqref{eq:HJexplicit}, however, the mixed term $p_1p_2\cos(\theta_1-\theta_2)$ and the factor
$(m_1+m_2\sin^2\Delta)$ in the denominator couple $\theta_1$ and $\theta_2$ irreducibly. No point transformation of the
angles is known that restores the Stäckel structure and, in accordance with the non-integrability result quoted above, none
exists in the analytic class. The Hamilton--Jacobi equation therefore has no closed-form complete integral in the
general case.

\subsection{Solvable limits}

\paragraph{(a) Small oscillations: separation in normal coordinates.}
Expand $H$ about the stable equilibrium $\bm\theta=0,\ \bm p=0$ to second order (see Section~\ref{sec:small}):
\begin{equation}
  H_2=\tfrac12\bm p^\top M_0^{-1}\bm p+\tfrac12\bm\theta^\top K\bm\theta,
  \qquad M_0=M(\bm0),\quad K=\mathrm{diag}(\mu g l_1,\ m_2g l_2).
\end{equation}
Let $A$ be the matrix of normalized modal vectors, $A^\top M_0A=\mathbb 1$, $A^\top KA=\mathrm{diag}(\omega_1^2,\omega_2^2)$
and define $\bm\theta=A\bm Q$ with conjugate momenta $\bm P=A^{\top}\bm p$ (hence $\bm P=\dot{\bm Q}$). Then
\begin{equation}
  H_2=\sum_{i=1}^{2}\tfrac12\big(P_i^2+\omega_i^2Q_i^2\big)
\end{equation}
is manifestly separable, and the Hamilton--Jacobi equation splits into two ordinary equations,
$\tfrac12(W_i')^2+\tfrac12\omega_i^2Q_i^2=\alpha_i$, with
\begin{equation}
  W(\bm Q,\bm\alpha)=\sum_{i=1}^{2}\int^{Q_i}\!\sqrt{2\alpha_i-\omega_i^2Q^2}\,\dd Q .
\end{equation}
The corresponding action--angle variables are
\begin{equation}
  J_i=\frac1{2\pi}\oint P_i\,\dd Q_i=\frac{\alpha_i}{\omega_i},\qquad
  \varphi_i=\frac{\partial W}{\partial J_i}=\omega_it+\varphi_i^{(0)},\qquad
  H_2=\omega_1J_1+\omega_2J_2 .
  \label{eq:AA}
\end{equation}
Phase-space trajectories in the linear approximation wind around invariant 2-tori $J_1,J_2=$ const.

\paragraph{(b) Anharmonic corrections: Birkhoff normal form.}
Retaining quartic terms in $H$ and applying a canonical transformation (via a generating function constructed order
by order in a Lie-series or Deprit scheme) one finds, away from resonances $k_1\omega_1+k_2\omega_2=0$,
\begin{equation}
  H=\omega_1J_1+\omega_2J_2+\tfrac12\big(a_{11}J_1^2+2a_{12}J_1J_2+a_{22}J_2^2\big)+O(J^3),
\end{equation}
where the $a_{ij}$ are rational functions of the masses and lengths. The frequencies $\Omega_i=\partial H/\partial J_i$
become amplitude-dependent. The normal form is generally only an asymptotic series; its divergence is the manifestation
of the destruction of tori by resonances (Section~\ref{sec:chaos}).

\paragraph{(c) Other solvable limits.}
When $m_2\to0$ the lower pendulum is a test particle, driven by the upper simple pendulum, and the system is a
one-degree-of-freedom pendulum plus a driven one: integrable in its autonomous part. When $l_2\to0$ it reduces to
a simple pendulum of mass $\mu$. When $m_1\to0$ the upper mass is a massless joint and the Hamiltonian is singular
at $\Delta=0,\pi$.

\section{Other Classical Formulations}
\label{sec:others}

\subsection{Hamilton's variational principle}

The motion makes the action $\mathcal S[\bm\theta]=\int_{t_0}^{t_1}L\,\dd t$ stationary among paths with fixed endpoints,
$\delta\mathcal S=0$. The Euler--Lagrange equations \eqref{eq:EL1}--\eqref{eq:EL2} are the Euler equations of this
variational problem. The equivalent phase-space form,
$\delta\int\big(\bm p\cdot\dd\bm\theta-H\,\dd t\big)=0$, with independent variations of $\bm\theta$ and $\bm p$,
reproduces Hamilton's canonical equations. The latter form is the starting point for symplectic (variational)
integrators.

\subsection{The Jacobi--Maupertuis principle and geodesic flow}

For fixed energy $E$, the trajectory (but not the time parametrization) is determined by the Maupertuis
(abbreviated-action) principle~\cite{arnold,abrahammarsden}
\begin{equation}
  \delta\int\bm p\cdot\dd\bm\theta=\delta\int\sqrt{2(E-V)\,M_{ij}\,\dd\theta^i\dd\theta^j}=0 .
\end{equation}
This is the statement that trajectories of energy $E$ are \emph{geodesics} of the Jacobi metric
\begin{equation}
  g^{(E)}_{ij}=2\,(E-V(\bm\theta))\,M_{ij}(\bm\theta)
\end{equation}
on the region of configuration space where $E>V$ (the Hill region). The geometry of this Riemannian two-manifold is
a direct handle on the dynamics. At low energies the Hill region is a small disc around the minimum of $V$, where the
Gaussian curvature of $g^{(E)}$ is positive and geodesics are stable. At higher energies the Hill region wraps around the
torus and includes regions of negative curvature, where nearby geodesics diverge exponentially. Negative curvature is
a sufficient (not necessary) mechanism for chaos.

\subsection{Gibbs--Appell equations}

Following Gibbs~\cite{gibbs} and Appell~\cite{appell}, let the \emph{acceleration energy} be
$\mathcal G=\tfrac12\sum_km_k\ddot{\bm r}_k^{\,2}$. Written in terms of the generalized accelerations $\ddot\theta_i$,
\begin{equation}
  \mathcal G=\tfrac12\ddot{\bm\theta}^{\top}M(\bm\theta)\ddot{\bm\theta}+\ddot{\bm\theta}\cdot\bm c(\bm\theta,\dot{\bm\theta})
  +\ (\text{terms independent of }\ddot{\bm\theta}),
\end{equation}
where $\bm c$ collects the velocity-dependent (Coriolis/centripetal) contributions. The Gibbs--Appell equations
\begin{equation}
  \frac{\partial\mathcal G}{\partial\ddot\theta_i}=Q_i,\qquad Q_i=-\frac{\partial V}{\partial\theta_i},
\end{equation}
are \emph{first-order in the accelerations}: they produce $M\ddot{\bm\theta}=-\bm c-\nabla V$, identical to
\eqref{eq:matrixform}. The formalism needs only the accelerations, never the constraint forces, and is particularly
efficient for nonholonomic systems; for holonomic systems such as the double pendulum it offers no savings over the
Lagrangian method but provides an independent check.

\subsection{Kane's method}
\label{sec:kane}

Kane's method~\cite{kane1961,kanewang,kanelevinson} is a vectorial formalism that, like the Lagrangian one,
eliminates workless constraint forces, but does so \emph{without} constructing the energy functions $T$ or $L$
and without differentiating them with respect to the coordinates. It is equivalent to Jourdain's principle of
virtual power and to the Gibbs--Appell equations, and it is the standard tool for generating the equations of
motion of multibody systems by computer algebra (for example in the Autolev and PyDy packages).

\paragraph{The recipe.} Consider $N$ particles $P_k$ with $n$ degrees of freedom.
\begin{enumerate}
  \item \textbf{Generalized speeds.} Choose $n$ independent quantities
        $u_r=\sum_sY_{rs}(\bm q,t)\,\dot q_s+z_r(\bm q,t)$, $r=1,\dots,n$, with $Y$ invertible. The simplest choice
        is $u_r=\dot q_r$. For a system with $m$ nonholonomic constraints the freedom in this choice is used to make
        the constraints solved identically, so that only $n-m$ independent speeds remain.
  \item \textbf{Partial velocities.} Write every velocity as linear in the speeds,
        $\bm v_k=\sum_r\bm v_k^{\,r}u_r+\bm v_k^{\,t}$. The coefficients
        $\bm v_k^{\,r}=\partial\bm v_k/\partial u_r$ are the \emph{partial velocities}. They span the directions along which
        the system can move without violating the constraints.
  \item \textbf{Accelerations.} Differentiate once in an inertial frame, $\bm a_k=\dot{\bm v}_k$.
  \item \textbf{Generalized active forces.}\quad $F_r=\sum_k\bm F_k\cdot\bm v_k^{\,r}$.
  \item \textbf{Generalized inertia forces.}\quad $F_r^*=-\sum_km_k\,\bm a_k\cdot\bm v_k^{\,r}$.
  \item \textbf{Kane's equations.}\quad $F_r+F_r^*=0$, $\quad r=1,\dots,n$.
\end{enumerate}
Ideal constraint forces never appear in $F_r$: they are orthogonal to the partial velocities (or, for internal
forces, cancel in pairs). In the present problem the rod forces are $-T_1\bm e_1$ on $m_1$, $+T_2\bm e_2$ on $m_1$ and
$-T_2\bm e_2$ on $m_2$, so their contribution to $F_r$ is
\[
  -T_1\bm e_1\cdot\bm v_1^{\,r}+T_2\bm e_2\cdot\big(\bm v_1^{\,r}-\bm v_2^{\,r}\big),
\]
which vanishes for $r=1,2$ because $\bm e_1\perp\bm t_1$ and $\bm e_2\perp\bm t_2$ (see the table below).

\paragraph{Application to the double pendulum.} Take $u_1=\dot\theta_1$, $u_2=\dot\theta_2$ and the vectors
$\bm e_i,\bm t_i$ of Section~\ref{sec:newton}.

\begin{center}
\renewcommand{\arraystretch}{1.3}
\begin{tabular}{@{}lll@{}}
\toprule
 & $m_1$ & $m_2$\\
\midrule
$\bm v_k$ & $l_1u_1\bm t_1$ & $l_1u_1\bm t_1+l_2u_2\bm t_2$\\
$\bm v_k^{\,1}$ & $l_1\bm t_1$ & $l_1\bm t_1$\\
$\bm v_k^{\,2}$ & $\bm 0$ & $l_2\bm t_2$\\
$\bm a_k$ & $l_1(\dot u_1\bm t_1-u_1^2\bm e_1)$ & $\bm a_1+l_2(\dot u_2\bm t_2-u_2^2\bm e_2)$\\
\bottomrule
\end{tabular}
\end{center}

With $\bm F_k=m_k\bm g$ and $\bm g\cdot\bm t_i=-g\sin\theta_i$, the generalized active forces are
\begin{equation}
  F_1=(m_1+m_2)\,\bm g\cdot l_1\bm t_1=-\mu g\,l_1\sin\theta_1,\qquad
  F_2=m_2\,\bm g\cdot l_2\bm t_2=-m_2g\,l_2\sin\theta_2 .
\end{equation}
Using $\bm t_1\!\cdot\!\bm t_2=\cos\Delta$, $\bm t_1\!\cdot\!\bm e_2=-\sin\Delta$, $\bm e_1\!\cdot\!\bm t_2=\sin\Delta$,
the accelerations projected on the partial velocities are $\bm a_1\!\cdot\!\bm t_1=l_1\dot u_1$,
$\bm a_2\!\cdot\!\bm t_1=l_1\dot u_1+l_2\dot u_2\cos\Delta+l_2u_2^2\sin\Delta$ and
$\bm a_2\!\cdot\!\bm t_2=l_1\dot u_1\cos\Delta-l_1u_1^2\sin\Delta+l_2\dot u_2$, so that
\begin{align}
  F_1^*&=-l_1\Big[\mu\,l_1\dot u_1+m_2l_2\dot u_2\cos\Delta+m_2l_2u_2^2\sin\Delta\Big],\\
  F_2^*&=-m_2l_2\Big[l_2\dot u_2+l_1\dot u_1\cos\Delta-l_1u_1^2\sin\Delta\Big].
\end{align}
Kane's equations $F_1+F_1^*=0$ and $F_2+F_2^*=0$ are, term by term, the Lagrange equations
\eqref{eq:EL1}--\eqref{eq:EL2}. No energy function was formed and no derivative of $T$ was taken.

\paragraph{Generalized mass matrix.} The coefficient of $\dot u_s$ in $-F_r^*$ is
$M_{rs}=\sum_km_k\,\bm v_k^{\,r}\!\cdot\bm v_k^{\,s}$. Here
$M_{11}=(m_1+m_2)l_1^2$, $M_{12}=m_2l_1l_2\,\bm t_1\!\cdot\!\bm t_2=m_2l_1l_2\cos\Delta$ and $M_{22}=m_2l_2^2$,
which is exactly the matrix \eqref{eq:M}. In general Kane's equations take the form
$M(\bm q)\,\dot{\bm u}=\bm f(\bm q,\bm u,t)$, which can be generated recursively for chains of rigid bodies.

\paragraph{Equivalence with Lagrange's equations.} If $u_r=\dot q_r$, then $\partial\bm v_k/\partial\dot q_r=\partial\bm r_k/\partial q_r$ and
\[
  \sum_km_k\,\bm a_k\!\cdot\!\frac{\partial\bm r_k}{\partial q_r}
  =\frac{\dd}{\dd t}\frac{\partial T}{\partial\dot q_r}-\frac{\partial T}{\partial q_r},
\]
so that $F_r^*=-\big(\frac{\dd}{\dd t}\frac{\partial T}{\partial\dot q_r}-\frac{\partial T}{\partial q_r}\big)$ while
$F_r=Q_r=-\partial V/\partial q_r$ for conservative forces. Kane's equations are then Lagrange's equations; their
advantages appear for nonholonomic systems and for other choices of speeds (for example body-fixed angular velocity
components of rigid bodies), where Lagrange's equations require quasi-coordinates or multipliers.

\paragraph{Recovering a constraint force: auxiliary speeds.} Because constraint forces do not enter $F_r$ they are
absent from the equations of motion. To bring a particular one into evidence, Kane introduces an
\emph{auxiliary} speed that releases the corresponding constraint. For the tension of the lower rod, let $u_3$ be a
fictitious rate of extension of rod~2,
\[
  \bm v_2=l_1u_1\bm t_1+l_2u_2\bm t_2+u_3\,\bm e_2,\qquad \bm v_2^{\,3}=\bm e_2,\quad\bm v_1^{\,3}=\bm 0 .
\]
The force $-T_2\bm e_2$ on $m_2$ now has a nonzero partial velocity and counts as an active force:
\[
  F_3=(m_2\bm g-T_2\bm e_2)\cdot\bm e_2=m_2g\cos\theta_2-T_2,\qquad
  F_3^*=-m_2\bm a_2\!\cdot\!\bm e_2=m_2\big(l_1\dot u_1\sin\Delta+l_1u_1^2\cos\Delta+l_2u_2^2\big),
\]
where $\bm a_2$ is evaluated after setting $u_3=\dot u_3=0$. The equation $F_3+F_3^*=0$ gives
\[
  T_2=m_2\big[g\cos\theta_2+l_1\dot u_1\sin\Delta+l_1u_1^2\cos\Delta+l_2u_2^2\big],
\]
identical to \eqref{eq:T2}. The tension of the upper rod is obtained in the same way from a fictitious speed
$u_4$ along $\bm e_1$ applied to both bobs. In practice one solves the two Kane equations for $\dot u_1,\dot u_2$ and
evaluates the tensions afterwards, exactly the economy of the Lagrangian approach combined with direct access to
the reaction forces.

\subsection{Summary of formulations}

\begin{table}[ht]
\centering
\begin{tabular}{@{}lllll@{}}
\toprule
Formulation & Variables & Order & Constraint forces & Typical use\\
\midrule
Newton & $\bm r_k$, $T_a$ & 2 & explicit & physical insight\\
Lagrange multipliers & $\bm r_k$, $\lambda_a$ & 2 (DAE) & explicit & multibody, forces\\
Lagrange (generalized) & $\theta_i$ & 2 & eliminated & derivations\\
Hamilton & $\theta_i,p_i$ & 1 & eliminated & symplectic structure\\
Hamilton--Jacobi & $S(\theta_i,t)$ & PDE & eliminated & integrability, AA variables\\
Jacobi--Maupertuis & $\theta_i$ at fixed $E$ & geodesic & eliminated & geometry of chaos\\
Gibbs--Appell, Kane & $\theta_i,\dot\theta_i$ & 2 & eliminated & multibody dynamics\\
\bottomrule
\end{tabular}
\caption{The formulations of this chapter, all equivalent for the double pendulum.}
\end{table}

\section{Small Oscillations and Normal Modes}
\label{sec:small}

\subsection{Linearization}

The stable equilibrium is $\theta_1=\theta_2=0$, where $V$ attains its global minimum $V_{\min}=-\mu g l_1-m_2g l_2$.
Expanding $T$ and $V$ to second order,
\begin{equation}
  L_2=\tfrac12\dot{\bm\theta}^\top M_0\dot{\bm\theta}-\tfrac12\bm\theta^\top K\bm\theta,
  \quad
  M_0=\begin{pmatrix}\mu\,l_1^2&m_2l_1l_2\\ m_2l_1l_2&m_2l_2^2\end{pmatrix},\
  K=\begin{pmatrix}\mu g\,l_1&0\\0&m_2g\,l_2\end{pmatrix}.
\end{equation}
The linear equations $M_0\ddot{\bm\theta}+K\bm\theta=0$ have solutions $\bm\theta=\bm u\,e^{i\omega t}$ provided
$\det(K-\omega^2M_0)=0$.

\subsection{Normal-mode frequencies}

Evaluating the determinant and dividing by $m_2l_1l_2$ gives the quadratic equation in $\omega^2$
\begin{equation}
  m_1l_1l_2\,\omega^4-\mu g\,(l_1+l_2)\,\omega^2+\mu g^2=0,
\end{equation}
whose roots are
\begin{equation}
  \omega_\pm^2=\frac{g}{2m_1l_1l_2}\Big[\mu\,(l_1+l_2)\pm\sqrt{\mu^2(l_1+l_2)^2-4m_1\mu\,l_1l_2}\Big].
  \label{eq:omega}
\end{equation}
The discriminant is positive since $\mu^2(l_1+l_2)^2-4m_1\mu l_1l_2\ge\mu(l_1+l_2)^2\,m_1\!\cdot\!0+\ldots>0$; more simply,
$\mu(l_1+l_2)^2-4m_1l_1l_2=\mu(l_1-l_2)^2+4m_2l_1l_2>0$.
The mode shapes follow from the first row of $(K-\omega^2M_0)\bm u=0$:
\begin{equation}
  \frac{u_2}{u_1}=\frac{\mu\,(g-\omega^2l_1)}{\omega^2m_2l_2}.
\end{equation}

\paragraph{Equal masses and lengths.} For $m_1=m_2=m$, $l_1=l_2=l$,
\begin{equation}
  \omega_\pm^2=\frac gl\,(2\pm\sqrt2),\qquad \frac{u_2}{u_1}=\pm\sqrt2 .
\end{equation}
The lower-frequency mode ($\omega_-^2=(2-\sqrt2)g/l$) is the in-phase swing with the lower bob having amplitude
$\sqrt2$ times that of the upper bob; the higher-frequency mode ($\omega_+^2=(2+\sqrt2)g/l$) is an anti-phase
motion in which the bobs move in opposite directions. The ratio $\omega_+/\omega_-=\sqrt{(2+\sqrt2)/(2-\sqrt2)}=1+\sqrt2$
is irrational, so a generic small-amplitude motion is quasi-periodic and never exactly repeats.

\section{Phase-Space Structure and Chaos}
\label{sec:chaos}

\subsection{Energy regimes}

The potential \eqref{eq:V} has four critical points on the torus, given by $\sin\theta_1=\sin\theta_2=0$:
\begin{center}
\begin{tabular}{@{}lll@{}}
\toprule
$(\theta_1,\theta_2)$ & $V$ & Type\\
\midrule
$(0,0)$ & $-\mu g l_1-m_2gl_2$ & minimum\\
$(0,\pi)$ & $-\mu g l_1+m_2gl_2$ & saddle\\
$(\pi,0)$ & $+\mu g l_1-m_2gl_2$ & saddle\\
$(\pi,\pi)$ & $+\mu g l_1+m_2gl_2$ & maximum\\
\bottomrule
\end{tabular}
\end{center}
These energies delimit the dynamical regimes. Below $V(0,\pi)$ the lower pendulum cannot pass over the top
while the upper rod hangs down; motion is confined to a neighbourhood of the minimum and is predominantly regular.
For energies above the lowest saddle the lower bob can flip over; the trajectory can then visit larger regions of the
torus, where the Jacobi metric has negative curvature, and most orbits become chaotic. For energies above
the maximum, $E>V(\pi,\pi)$, the whole torus is accessible, and both bobs can rotate continuously. Detailed numerical and experimental maps of the regular and chaotic regions, including flip-time diagrams, can be found in~\cite{shinbrot,levien,stachowiak}.

\subsection{Poincar\'e sections}

For fixed $E$ the motion lies on a three-dimensional energy surface $\{H=E\}$. A convenient two-dimensional Poincar\'e
section is $\Sigma=\{\theta_1=0,\ \dot\theta_1>0\}$ with coordinates $(\theta_2,p_2)$; $p_1$ is then fixed by
$H=E$. Successive piercings of $\Sigma$ define an area-preserving map $(\theta_2,p_2)_n\to(\theta_2,p_2)_{n+1}$.
\begin{itemize}
  \item At low $E$ the points lie on closed invariant curves: the cross-sections of KAM tori.
  \item With increasing $E$, curves with rational frequency ratios $\Omega_1/\Omega_2=r/s$ break up into chains of
        $s$ islands (Poincar\'e--Birkhoff theorem), separated by thin chaotic layers around hyperbolic points.
  \item At still higher $E$ the last KAM tori are destroyed and the chaotic sea occupies most of the section, with
        residual stable islands.
\end{itemize}

\subsection{KAM theorem and Lyapunov exponents}

The Kolmogorov--Arnold--Moser theorem~\cite{kolmogorov,arnoldkam,moser} states that, for a small enough perturbation of an integrable Hamiltonian
$H_0(\bm J)$ and a sufficiently non-degenerate $H_0$, the tori with sufficiently irrational (Diophantine) frequency
vectors persist, slightly deformed. The double pendulum at low energy is a perturbation of the integrable
$H_2$ of \eqref{eq:AA} and so the theorem governs the picture near the minimum, subject to the usual non-degeneracy
condition, $\det(\partial^2H_0/\partial J_i\partial J_j)\ne0$, on the quartic coefficients $a_{ij}$.

Chaos is quantified by the largest Lyapunov exponent~\cite{benettin}
\begin{equation}
  \lambda_{\max}=\lim_{t\to\infty}\lim_{\|\delta z(0)\|\to0}\frac1t\ln\frac{\|\delta z(t)\|}{\|\delta z(0)\|},
\end{equation}
where $\delta z$ obeys the variational equation $\dot{\delta z}=J(z(t))\,\delta z$, $J=\partial\dot z/\partial z$.
Because the flow is Hamiltonian, the Lyapunov spectrum of this four-dimensional system is symmetric about zero,
$\{\lambda_{\max},0,0,-\lambda_{\max}\}$, where the two zero exponents correspond to the flow direction and to the
energy-conservation direction. A positive $\lambda_{\max}$ implies
sensitive dependence on initial conditions: two initial conditions that differ by $10^{-10}$ in an angle typically
separate to $O(1)$ over a time of $\sim\ln(10^{10})/\lambda_{\max}$.

\section{Numerical Integration}
\label{sec:numerics}

\subsection{Choice of method}

Since a closed-form solution is not available, numerical integration is the practical tool. Three remarks.

\begin{enumerate}
  \item The accelerations can be computed in the Lagrangian form by solving the $2\times2$ linear system
        $M\ddot{\bm\theta}=\bm f$ at each step, avoiding the cancellation-prone closed forms
        \eqref{eq:acc1}--\eqref{eq:acc2} close to $\Delta=0$.
  \item A standard explicit Runge--Kutta method of fourth order (RK4~\cite{hnw}) with a small step is adequate for visualizing
        the trajectory but is not symplectic: the energy drifts secularly over long integrations.
  \item The Hamiltonian \eqref{eq:H} is \emph{not separable} ($H\ne T(\bm p)+V(\bm\theta)$ because $T$ depends on
        $\bm\theta$), so the explicit leapfrog (St\"ormer--Verlet) scheme is not available. Symplectic alternatives are
        the implicit midpoint rule (second order, symplectic for any $H$), the generalized implicit St\"ormer--Verlet
        scheme of Hairer--Lubich--Wanner~\cite{ghw}, and higher-order Gauss--Legendre collocation methods. Variational integrators~\cite{marsdenwest}
        derived by discretizing the action give a closely related family.
\end{enumerate}

\subsection{A reference implementation in C}

The following program integrates the Lagrangian equations \eqref{eq:EL1}--\eqref{eq:EL2} with RK4 and prints the
relative energy error, the cheapest check of the quality of the integration. A Fortran version, which also
contains the symplectic integrators, is given in Section~\ref{sec:fortran}.

\begin{lstlisting}
#include <math.h>
#include <stdio.h>

typedef struct { double m1, m2, l1, l2, g; } Par;

/* y = { th1, th2, w1, w2 },  w = d(theta)/dt */
static void rhs(const double y[4], double dy[4], const Par *p)
{
    double d = y[0] - y[1], c = cos(d), s = sin(d);
    double mu = p->m1 + p->m2;
    double w1 = y[2], w2 = y[3];

    /* A * (th1'', th2'')^T = b */
    double a11 = mu * p->l1,  a12 = p->m2 * p->l2 * c;
    double a21 = p->l1 * c,   a22 = p->l2;
    double b1  = -p->m2 * p->l2 * w2 * w2 * s - mu * p->g * sin(y[0]);
    double b2  =  p->l1 * w1 * w1 * s - p->g * sin(y[1]);
    double det = a11 * a22 - a12 * a21;

    dy[0] = w1;
    dy[1] = w2;
    dy[2] = (b1 * a22 - a12 * b2) / det;
    dy[3] = (a11 * b2 - a21 * b1) / det;
}

static double energy(const double y[4], const Par *p)
{
    double d = y[0] - y[1], mu = p->m1 + p->m2;
    double T = 0.5 * mu * p->l1 * p->l1 * y[2] * y[2]
             + 0.5 * p->m2 * p->l2 * p->l2 * y[3] * y[3]
             + p->m2 * p->l1 * p->l2 * y[2] * y[3] * cos(d);
    double V = -mu * p->g * p->l1 * cos(y[0])
               - p->m2 * p->g * p->l2 * cos(y[1]);
    return T + V;
}

static void rk4(double y[4], double h, const Par *p)
{
    double k1[4], k2[4], k3[4], k4[4], t[4];
    int i;
    rhs(y, k1, p);
    for (i = 0; i < 4; i++) t[i] = y[i] + 0.5 * h * k1[i];
    rhs(t, k2, p);
    for (i = 0; i < 4; i++) t[i] = y[i] + 0.5 * h * k2[i];
    rhs(t, k3, p);
    for (i = 0; i < 4; i++) t[i] = y[i] + h * k3[i];
    rhs(t, k4, p);
    for (i = 0; i < 4; i++)
        y[i] += h * (k1[i] + 2*k2[i] + 2*k3[i] + k4[i]) / 6.0;
}

int main(void)
{
    Par p = { 1.0, 1.0, 1.0, 1.0, 9.81 };
    double y[4] = { 2.0, 2.5, 0.0, 0.0 };   /* released from rest */
    double h = 1e-3, E0 = energy(y, &p);
    long n;

    for (n = 0; n <= 100000; n++) {          /* t = 100 */
        if (n % 1000 == 0)
            printf("%8.3f %12.6f %12.6f %12.3e\n",
                   n * h, y[0], y[1], (energy(y, &p) - E0) / fabs(E0));
        rk4(y, h, &p);
    }
    return 0;
}
\end{lstlisting}

\subsection{Symplectic integrators for a non-separable Hamiltonian}
\label{sec:symplectic}

A one-step map $\Phi_h:z_n\mapsto z_{n+1}$ is \emph{symplectic} if its Jacobian satisfies
$(\Phi_h')^{\top}\mathbb J\,\Phi_h'=\mathbb J$, i.e.\ if it preserves the form $\omega$ exactly, as the true flow
does~\cite{ghw,sanzserna,leimkuhler}. Backward error analysis shows that a symplectic method of order $p$ is the exact
flow of a modified Hamiltonian $\tilde H=H+h^pH_{p+1}+h^{p+1}H_{p+2}+\dots$ (an asymptotic series). Consequently $H$ is
conserved up to a bounded, oscillating error of size $O(h^p)$ over times that are exponentially long in $1/h$, in
contrast with the secular drift of a generic non-symplectic method. The same argument shows that the
implementation must solve the implicit equations to (near) machine precision; a loose iteration tolerance destroys
the symplectic character and with it the long-time behaviour.

Write $f(z)=\big(\partial H/\partial\bm p,\,-\partial H/\partial\bm\theta\big)$ for the Hamiltonian vector field, with
$\partial H/\partial\bm p=M^{-1}\bm p=\dot{\bm\theta}$ and $\partial H/\partial\bm\theta=-\partial L/\partial\bm\theta$ as
in Section~\ref{sec:hamilton}. Three symplectic schemes are used below.

\paragraph{(a) Implicit midpoint rule.} 
\begin{equation}
  z_{n+1}=z_n+h\,f\!\Big(\frac{z_n+z_{n+1}}2\Big).
\end{equation}
It is symplectic for \emph{every} $H$, symmetric (time-reversible), of order~2, and conserves all quadratic first
integrals.

\paragraph{(b) Generalized St\"ormer--Verlet.} For a non-separable $H(\bm\theta,\bm p)$ the explicit leapfrog is not
available. Its symplectic generalization~\cite{ghw} is
\begin{align}
  \bm p_{n+1/2}&=\bm p_n-\tfrac h2\,\partial_{\bm\theta}H(\bm\theta_n,\bm p_{n+1/2}),\\
  \bm\theta_{n+1}&=\bm\theta_n+\tfrac h2\big[\partial_{\bm p}H(\bm\theta_n,\bm p_{n+1/2})
                   +\partial_{\bm p}H(\bm\theta_{n+1},\bm p_{n+1/2})\big],\\
  \bm p_{n+1}&=\bm p_{n+1/2}-\tfrac h2\,\partial_{\bm\theta}H(\bm\theta_{n+1},\bm p_{n+1/2}).
\end{align}
It is the composition of two symplectic Euler steps of length $h/2$, hence symplectic, symmetric and of order~2. The
first two equations are implicit in $\bm p_{n+1/2}$ and $\bm\theta_{n+1}$ respectively and are solved by fixed-point
iteration; the third is explicit. For $H=T(\bm p)+V(\bm\theta)$ the scheme reduces to the usual explicit leapfrog.

\paragraph{(c) Fourth-order composition.} If $\Psi_h$ is a symmetric method of order 2, then the triple jump
\begin{equation}
  \Psi^{[4]}_h=\Psi_{\gamma_1h}\circ\Psi_{\gamma_0h}\circ\Psi_{\gamma_1h},\qquad
  \gamma_1=\frac1{2-2^{1/3}},\quad\gamma_0=1-2\gamma_1=-\frac{2^{1/3}}{2-2^{1/3}},
\end{equation}
has order 4~\cite{yoshida} (the conditions are $2\gamma_1+\gamma_0=1$ and $2\gamma_1^3+\gamma_0^3=0$) and remains
symplectic and symmetric, as compositions of symplectic maps are. Because $\gamma_0<0$ the middle step runs backwards
in time. The cost is three Verlet steps per step. Here $\Psi_h$ is the scheme (b).

\paragraph{Solving the implicit equations.} For a vector field with Lipschitz constant $\mathcal L$ the fixed-point
iterations converge if $h\mathcal L<1$ (for the double pendulum $\mathcal L$ grows with the angular velocities and
with $1/\det M$). If a larger step is needed, replace the fixed-point iteration by Newton's method with the Hessian
of $H$ (analytic or by finite differences), at the cost of a $4\times4$ linear solve per iteration.

\subsection{A reference implementation in Fortran}
\label{sec:fortran}

The Fortran~2008 program below collects all four integrators in a module \texttt{dp\_mod}. The Lagrangian state is
$y=(\theta_1,\theta_2,\dot\theta_1,\dot\theta_2)$ for RK4 and the canonical state is
$z=(\theta_1,\theta_2,p_1,p_2)$ for the symplectic methods; \texttt{vel\_to\_mom} and \texttt{mom\_to\_vel} convert
between them via the mass matrix \eqref{eq:M} and its inverse \eqref{eq:thd1}--\eqref{eq:thd2}. The routines
\texttt{h\_p} and \texttt{h\_th} are the analytic partial derivatives of \eqref{eq:H} (the latter using
$\partial H/\partial\bm\theta=-\partial L/\partial\bm\theta$), and \texttt{vfield} assembles the vector field $f$.
The driver integrates two initial conditions with the same step $h=10^{-2}$ for $t\in[0,1000]$ and records the
relative energy error of every method. Compile with
\begin{center}
\texttt{gfortran -O2 -std=f2008 -Wall double\_pendulum.f90 -o dp}
\end{center}

\begin{lstlisting}[language={[95]Fortran},basicstyle=\ttfamily\scriptsize]
! double_pendulum.f90
! Planar double pendulum: RK4 (Lagrangian form) and three symplectic
! integrators for the non-separable Hamiltonian H(theta,p):
!   - implicit midpoint rule            (order 2)
!   - generalized Stormer-Verlet        (order 2)
!   - Yoshida triple-jump composition   (order 4)
!
! Build:  gfortran -O2 -std=f2008 -Wall double_pendulum.f90 -o dp
!
! Conventions: theta_i measured from the downward vertical,
!   Lagrangian state  y = (th1, th2, w1, w2),   w = d(theta)/dt
!   canonical state   z = (th1, th2, p1, p2)
module dp_mod
  use, intrinsic :: iso_fortran_env, only: wp => real64
  implicit none
  private
  public :: wp, par_t, rhs, energy_v, ham, vel_to_mom, mom_to_vel
  public :: rk4_step, midpoint_step, sv_step, yoshida4_step

  type :: par_t
     real(wp) :: m1 = 1.0_wp, m2 = 1.0_wp
     real(wp) :: l1 = 1.0_wp, l2 = 1.0_wp
     real(wp) :: g  = 9.81_wp
  end type par_t

  integer,  parameter :: maxit = 200
  real(wp), parameter :: tol   = 1.0e-13_wp

contains

  !---------------------------------------------------------------
  ! Lagrangian right-hand side: solve A*(th1'',th2'')^T = b
  !---------------------------------------------------------------
  subroutine rhs(p, y, dy)
    type(par_t), intent(in)  :: p
    real(wp),    intent(in)  :: y(4)
    real(wp),    intent(out) :: dy(4)
    real(wp) :: d, c, s, mu, a11, a12, a21, a22, b1, b2, det

    d  = y(1) - y(2);  c = cos(d);  s = sin(d);  mu = p%m1 + p%m2
    a11 = mu*p%l1;  a12 = p%m2*p%l2*c
    a21 = p%l1*c;   a22 = p%l2
    b1  = -p%m2*p%l2*y(4)**2*s - mu*p%g*sin(y(1))
    b2  =  p%l1*y(3)**2*s      - p%g*sin(y(2))
    det = a11*a22 - a12*a21

    dy(1:2) = y(3:4)
    dy(3)   = (b1*a22 - a12*b2)/det
    dy(4)   = (a11*b2 - a21*b1)/det
  end subroutine rhs

  !---------------------------------------------------------------
  ! Total energy from Lagrangian variables (th, w)
  !---------------------------------------------------------------
  function energy_v(p, y) result(e)
    type(par_t), intent(in) :: p
    real(wp),    intent(in) :: y(4)
    real(wp) :: e, mu, t, v
    mu = p%m1 + p%m2
    t = 0.5_wp*mu*p%l1**2*y(3)**2 + 0.5_wp*p%m2*p%l2**2*y(4)**2        &
        + p%m2*p%l1*p%l2*y(3)*y(4)*cos(y(1) - y(2))
    v = -mu*p%g*p%l1*cos(y(1)) - p%m2*p%g*p%l2*cos(y(2))
    e = t + v
  end function energy_v

  !---------------------------------------------------------------
  ! Momenta <-> velocities (inverse mass matrix, eqs. thd1, thd2)
  !---------------------------------------------------------------
  subroutine vel_to_mom(p, th, w, mom)
    type(par_t), intent(in)  :: p
    real(wp),    intent(in)  :: th(2), w(2)
    real(wp),    intent(out) :: mom(2)
    real(wp) :: c, mu
    c = cos(th(1) - th(2));  mu = p%m1 + p%m2
    mom(1) = mu*p%l1**2*w(1) + p%m2*p%l1*p%l2*c*w(2)
    mom(2) = p%m2*p%l2**2*w(2) + p%m2*p%l1*p%l2*c*w(1)
  end subroutine vel_to_mom

  subroutine mom_to_vel(p, th, mom, w)
    type(par_t), intent(in)  :: p
    real(wp),    intent(in)  :: th(2), mom(2)
    real(wp),    intent(out) :: w(2)
    real(wp) :: c, s, mu, den
    c = cos(th(1) - th(2));  s = sin(th(1) - th(2));  mu = p%m1 + p%m2
    den  = p%m1 + p%m2*s*s
    w(1) = (p%l2*mom(1) - p%l1*c*mom(2)) / (p%l1**2*p%l2*den)
    w(2) = (mu*p%l1*mom(2) - p%m2*p%l2*c*mom(1)) / (p%m2*p%l1*p%l2**2*den)
  end subroutine mom_to_vel

  !---------------------------------------------------------------
  ! Hamiltonian, eq. (H), and its partial derivatives
  !---------------------------------------------------------------
  function ham(p, th, mom) result(h)
    type(par_t), intent(in) :: p
    real(wp),    intent(in) :: th(2), mom(2)
    real(wp) :: h, w(2), mu
    mu = p%m1 + p%m2
    call mom_to_vel(p, th, mom, w)
    h = 0.5_wp*(mom(1)*w(1) + mom(2)*w(2))                              &
        - mu*p%g*p%l1*cos(th(1)) - p%m2*p%g*p%l2*cos(th(2))
  end function ham

  subroutine h_p(p, th, mom, hp)             ! dH/dp = M^{-1} p = theta'
    type(par_t), intent(in)  :: p
    real(wp),    intent(in)  :: th(2), mom(2)
    real(wp),    intent(out) :: hp(2)
    call mom_to_vel(p, th, mom, hp)
  end subroutine h_p

  subroutine h_th(p, th, mom, hq)            ! dH/dtheta at fixed p
    type(par_t), intent(in)  :: p
    real(wp),    intent(in)  :: th(2), mom(2)
    real(wp),    intent(out) :: hq(2)
    real(wp) :: w(2), s, mu, k
    mu = p%m1 + p%m2
    call mom_to_vel(p, th, mom, w)
    s = sin(th(1) - th(2))
    k = p%m2*p%l1*p%l2*w(1)*w(2)*s
    hq(1) =  k + mu*p%g*p%l1*sin(th(1))
    hq(2) = -k + p%m2*p%g*p%l2*sin(th(2))
  end subroutine h_th

  subroutine vfield(p, z, f)                 ! f = (dH/dp, -dH/dtheta)
    type(par_t), intent(in)  :: p
    real(wp),    intent(in)  :: z(4)
    real(wp),    intent(out) :: f(4)
    real(wp) :: hq(2)
    call h_p (p, z(1:2), z(3:4), f(1:2))
    call h_th(p, z(1:2), z(3:4), hq)
    f(3:4) = -hq
  end subroutine vfield

  !---------------------------------------------------------------
  ! Classical RK4 on the Lagrangian first-order system
  !---------------------------------------------------------------
  subroutine rk4_step(p, y, h)
    type(par_t), intent(in)    :: p
    real(wp),    intent(inout) :: y(4)
    real(wp),    intent(in)    :: h
    real(wp) :: k1(4), k2(4), k3(4), k4(4)
    call rhs(p, y,              k1)
    call rhs(p, y + 0.5_wp*h*k1, k2)
    call rhs(p, y + 0.5_wp*h*k2, k3)
    call rhs(p, y +        h*k3, k4)
    y = y + h*(k1 + 2.0_wp*k2 + 2.0_wp*k3 + k4)/6.0_wp
  end subroutine rk4_step

  !---------------------------------------------------------------
  ! Implicit midpoint rule  z1 = z0 + h f((z0+z1)/2)
  ! (symplectic and symmetric for any H); fixed-point iteration,
  ! convergent for h*|Hessian H| < 1.  Use Newton for larger h.
  !---------------------------------------------------------------
  subroutine midpoint_step(p, z, h, ok)
    type(par_t), intent(in)    :: p
    real(wp),    intent(inout) :: z(4)
    real(wp),    intent(in)    :: h
    logical,     intent(out)   :: ok
    real(wp) :: f(4), zn(4), zt(4), err
    integer  :: it

    call vfield(p, z, f)
    zn = z + h*f                              ! explicit Euler predictor
    ok = .false.
    do it = 1, maxit
       call vfield(p, 0.5_wp*(z + zn), f)
       zt  = z + h*f
       err = maxval(abs(zt - zn))
       zn  = zt
       if (err < tol*(1.0_wp + maxval(abs(zn)))) then
          ok = .true.;  exit
       end if
    end do
    z = zn
  end subroutine midpoint_step

  !---------------------------------------------------------------
  ! Generalized (implicit) Stormer-Verlet, Hairer-Lubich-Wanner:
  !   p_{1/2} = p_0 - h/2 H_th(q_0, p_{1/2})
  !   q_1     = q_0 + h/2 [H_p(q_0,p_{1/2}) + H_p(q_1,p_{1/2})]
  !   p_1     = p_{1/2} - h/2 H_th(q_1, p_{1/2})
  !---------------------------------------------------------------
  subroutine sv_step(p, z, h, ok)
    type(par_t), intent(in)    :: p
    real(wp),    intent(inout) :: z(4)
    real(wp),    intent(in)    :: h
    logical,     intent(out)   :: ok
    real(wp) :: q0(2), p0(2), ph(2), pt(2), q1(2), qt(2)
    real(wp) :: hq(2), hp0(2), hp1(2), err
    integer  :: it
    logical  :: ok1, ok2

    q0 = z(1:2);  p0 = z(3:4)

    ! (1) implicit half step in p
    ph = p0;  ok1 = .false.
    do it = 1, maxit
       call h_th(p, q0, ph, hq)
       pt  = p0 - 0.5_wp*h*hq
       err = maxval(abs(pt - ph));  ph = pt
       if (err < tol*(1.0_wp + maxval(abs(ph)))) then
          ok1 = .true.;  exit
       end if
    end do

    ! (2) implicit full step in q
    call h_p(p, q0, ph, hp0)
    q1 = q0 + h*hp0;  ok2 = .false.
    do it = 1, maxit
       call h_p(p, q1, ph, hp1)
       qt  = q0 + 0.5_wp*h*(hp0 + hp1)
       err = maxval(abs(qt - q1));  q1 = qt
       if (err < tol*(1.0_wp + maxval(abs(q1)))) then
          ok2 = .true.;  exit
       end if
    end do

    ! (3) explicit half step in p
    call h_th(p, q1, ph, hq)
    z(1:2) = q1
    z(3:4) = ph - 0.5_wp*h*hq
    ok = ok1 .and. ok2
  end subroutine sv_step

  !---------------------------------------------------------------
  ! Order-4 composition (Yoshida 1990 / Suzuki triple jump):
  !   Psi4(h) = Psi2(g1 h) o Psi2(g0 h) o Psi2(g1 h)
  !---------------------------------------------------------------
  subroutine yoshida4_step(p, z, h, ok)
    type(par_t), intent(in)    :: p
    real(wp),    intent(inout) :: z(4)
    real(wp),    intent(in)    :: h
    logical,     intent(out)   :: ok
    real(wp), parameter :: g1 = 1.0_wp/(2.0_wp - 2.0_wp**(1.0_wp/3.0_wp))
    real(wp), parameter :: g0 = 1.0_wp - 2.0_wp*g1
    logical :: o1, o2, o3
    call sv_step(p, z, g1*h, o1)
    call sv_step(p, z, g0*h, o2)
    call sv_step(p, z, g1*h, o3)
    ok = o1 .and. o2 .and. o3
  end subroutine yoshida4_step

end module dp_mod

!=====================================================================
program dp_compare
  use dp_mod
  implicit none
  type(par_t) :: p
  real(wp) :: th0(2), w0(2), mom0(2), h, tfin, e0, emax(4)
  real(wp) :: yr(4), zm(4), zs(4), zy(4)
  integer  :: icase, n, nsteps, nprint
  logical  :: ok

  h      = 1.0e-2_wp
  tfin   = 1000.0_wp
  nsteps = nint(tfin/h)
  nprint = nsteps/10

  do icase = 1, 2
     if (icase == 1) then                      ! low energy: regular motion
        th0 = [0.4_wp, 0.2_wp]
        print '(a)', '# case 1: th0 = (0.4, 0.2), at rest (regular regime)'
     else                                      ! high energy: chaotic
        th0 = [2.0_wp, 2.5_wp]
        print '(a)', '# case 2: th0 = (2.0, 2.5), at rest (chaotic regime)'
     end if
     w0 = 0.0_wp
     call vel_to_mom(p, th0, w0, mom0)

     yr = [th0, w0]
     zm = [th0, mom0];  zs = zm;  zy = zm
     e0 = ham(p, th0, mom0)
     emax = 0.0_wp
     print '(a)', '#    t        RK4          midpoint     Stormer-Verlet  Yoshida4'

     do n = 1, nsteps
        call rk4_step(p, yr, h)
        call midpoint_step(p, zm, h, ok);   if (.not. ok) print *, 'midpoint: no conv.'
        call sv_step(p, zs, h, ok);         if (.not. ok) print *, 'SV: no conv.'
        call yoshida4_step(p, zy, h, ok);   if (.not. ok) print *, 'Y4: no conv.'

        emax(1) = max(emax(1), abs(energy_v(p, yr) - e0))
        emax(2) = max(emax(2), abs(ham(p, zm(1:2), zm(3:4)) - e0))
        emax(3) = max(emax(3), abs(ham(p, zs(1:2), zs(3:4)) - e0))
        emax(4) = max(emax(4), abs(ham(p, zy(1:2), zy(3:4)) - e0))

        if (mod(n, nprint) == 0) then
           print '(f8.1,4es14.3)', n*h,                                 &
              (energy_v(p, yr) - e0)/abs(e0),                           &
              (ham(p, zm(1:2), zm(3:4)) - e0)/abs(e0),                  &
              (ham(p, zs(1:2), zs(3:4)) - e0)/abs(e0),                  &
              (ham(p, zy(1:2), zy(3:4)) - e0)/abs(e0)
        end if
     end do
     print '(a,4es14.3)', '# max |dE/E0| :', emax/abs(e0)
     print *
  end do
end program dp_compare
\end{lstlisting}

\subsection{Verification and results}
\label{sec:results}

\paragraph{Order of convergence.} For the initial condition $\bm\theta(0)=(2.0,\,2.5)$ at rest
(equal masses and lengths, $g=9.81$) the solution at $t=2$ was computed with each method for a sequence of step
sizes and compared with an RK4 reference obtained with $h=10^{-5}$. The maximum angle error is listed in
Table~\ref{tab:order}. The observed orders agree with the theory: 2 for midpoint and Verlet, 4 for the composition
and for RK4 (the apparent order of RK4 is still approaching 4 at the largest steps).

\begin{table}[ht]
\centering
\begin{tabular}{@{}lcccc@{}}
\toprule
$h$ & RK4 & midpoint & St\"ormer--Verlet & Yoshida-4\\
\midrule
$2.0\times10^{-2}$ & $2.2\times10^{-4}$ & $4.8\times10^{-3}$ & $3.0\times10^{-3}$ & $1.0\times10^{-3}$\\
$1.0\times10^{-2}$ & $8.8\times10^{-6}$ & $1.2\times10^{-3}$ & $7.5\times10^{-4}$ & $7.2\times10^{-5}$\\
$5.0\times10^{-3}$ & $3.7\times10^{-7}$ & $3.0\times10^{-4}$ & $1.9\times10^{-4}$ & $4.7\times10^{-6}$\\
$2.5\times10^{-3}$ & $1.7\times10^{-8}$ & $7.5\times10^{-5}$ & $4.7\times10^{-5}$ & $2.9\times10^{-7}$\\
$1.25\times10^{-3}$ & $8.9\times10^{-10}$ & $1.9\times10^{-5}$ & $1.2\times10^{-5}$ & $1.8\times10^{-8}$\\
\midrule
observed order & 4.3 & 2.00 & 2.00 & 4.00\\
\bottomrule
\end{tabular}
\caption{Maximum angle error at $t=2$ versus step size, $\bm\theta(0)=(2.0,2.5)$.}
\label{tab:order}
\end{table}

\paragraph{Energy conservation.} Table~\ref{tab:energy} shows the relative energy error $(E-E_0)/|E_0|$ at
$t=1000$ and its maximum modulus over the run, for the regular case $\bm\theta(0)=(0.4,0.2)$ and the chaotic case
$\bm\theta(0)=(2.0,2.5)$, both at rest, with $h=10^{-2}$.

\begin{table}[ht]
\centering
\begin{tabular}{@{}llcccc@{}}
\toprule
case & quantity & RK4 & midpoint & St\"ormer--Verlet & Yoshida-4\\
\midrule
regular & error at $t=1000$ & $-5.5\times10^{-7}$ & $6.4\times10^{-8}$ & $-6.4\times10^{-6}$ & $2.2\times10^{-9}$\\
        & $\max|\Delta E/E_0|$ & $5.5\times10^{-7}$ & $8.4\times10^{-7}$ & $1.7\times10^{-5}$ & $8.5\times10^{-9}$\\
chaotic & error at $t=1000$ & $-1.5\times10^{-2}$ & $3.4\times10^{-4}$ & $5.7\times10^{-5}$ & $1.4\times10^{-8}$\\
        & $\max|\Delta E/E_0|$ & $1.5\times10^{-2}$ & $1.2\times10^{-2}$ & $2.4\times10^{-2}$ & $3.7\times10^{-3}$\\
\bottomrule
\end{tabular}
\caption{Relative energy error for $h=10^{-2}$, $t\in[0,1000]$.}
\label{tab:energy}
\end{table}

The qualitative difference is in the \emph{time dependence}, not in the size at a given instant. The RK4 error grows
monotonically (linearly in the regular case, and to $1.5\%$ in the chaotic one), while the symplectic errors oscillate
without trend. In the regular case the second-order symplectic schemes have bounded errors that are comparable to, or larger than,
the RK4 error at $t=1000$, but the RK4 error keeps growing linearly and would exceed them in longer runs; only the
fourth-order composition is better than RK4 at every time in this test. In the chaotic case the large
\emph{maximum} for the symplectic methods is a transient excursion, not a drift: the sampled values in the program
output fluctuate in sign and do not grow. Such excursions are associated with phases of fast motion, where the step
$h=10^{-2}$ is coarse for a second-order method. The size of the oscillation decreases as $h^p$, so the step should be
chosen according to the fastest time scale visited by the trajectory rather than the average one. As always in the chaotic regime, no integrator reproduces an individual
orbit beyond a few Lyapunov times; the symplectic property guarantees the quality of the \emph{statistics}.

\section{Summary}

\begin{itemize}
  \item The double pendulum has two degrees of freedom and the kinetic energy
        \eqref{eq:T} with the configuration-dependent mass matrix \eqref{eq:M}.
  \item The Newtonian, Lagrangian, and Hamiltonian formulations are equivalent; they differ in whether constraint
        forces appear, the order of the differential equations, and the geometric structure that is made manifest.
  \item The Hamiltonian \eqref{eq:H} is non-separable; the Hamilton--Jacobi equation is solvable only in the small-oscillation
        limit, where action--angle variables \eqref{eq:AA} exist.
  \item Energy is the only known integral. At low energy KAM theory ensures predominantly regular motion; above the saddle
        energies the phase space becomes largely chaotic.
  \item Kane's method gives the same equations as Lagrange's without forming $T$ or $L$, and yields the constraint
        forces directly through auxiliary speeds (Section~\ref{sec:kane}).
  \item Numerical integration must respect the structure of the problem: symplectic and reversible schemes for
        long-time statistics, and energy monitoring as a minimal check of every run.
\end{itemize}

\section*{Exercises}
\addcontentsline{toc}{section}{Exercises}

\begin{exercise}
Starting from \eqref{eq:EL1}--\eqref{eq:EL2}, verify the explicit accelerations \eqref{eq:acc1}--\eqref{eq:acc2}.
\end{exercise}

\begin{exercise}
Show that $\det M=m_2l_1^2l_2^2(m_1+m_2\sin^2\Delta)$ and compute the inverse mass matrix. Hence derive the
Hamiltonian \eqref{eq:H} and check that its Legendre transform returns the Lagrangian.
\end{exercise}

\begin{exercise}
Verify the claimed identity $\partial H/\partial\theta_i|_{\bm p}=-\partial L/\partial\theta_i|_{\dot{\bm\theta}}$ for
a Lagrangian $L=\tfrac12\dot q^\top M(q)\dot q-V(q)$ and use it to derive $\dot p_1$, $\dot p_2$ in
Section~\ref{sec:hamilton}.
\end{exercise}

\begin{exercise}
Compute $\pb{p_1+p_2}{H}$ and show that it equals $-\mu g l_1\sin\theta_1-m_2gl_2\sin\theta_2$ (the total gravitational torque about $O$).
Interpret the result, and show that for $g=0$ the quantity $p_1+p_2$ is a second integral.
\end{exercise}

\begin{exercise}
Find the normal-mode frequencies \eqref{eq:omega} in the limits $m_1\gg m_2$ and $m_1\ll m_2$, and for
$l_2\to0$. Interpret each limit physically.
\end{exercise}

\begin{exercise}
Derive the quartic terms of $H$ about the equilibrium and compute the coefficients $a_{11},a_{12},a_{22}$ of the Birkhoff normal
form for $m_1=m_2$, $l_1=l_2$. Identify the lowest-order resonances $\omega_+/\omega_-=r/s$ with small $r,s$ that can be
reached by tuning $m_2/m_1$ and $l_2/l_1$.
\end{exercise}

\begin{exercise}
Compute the Gaussian curvature of the Jacobi metric $g^{(E)}_{ij}=2(E-V)M_{ij}$ on the torus and locate the regions
where it is negative for $m_1=m_2$, $l_1=l_2$ and several values of $E$. Compare with the Poincar\'e section computed
numerically.
\end{exercise}

\begin{exercise}
Using the program of Section~\ref{sec:numerics}, compute the largest Lyapunov exponent by integrating the variational
equations alongside the trajectory (or by the two-trajectory method with periodic renormalization). Plot $\lambda_{\max}$
as a function of $E$ for $m_1=m_2$, $l_1=l_2$, and compare with the energies of Section~\ref{sec:chaos}.
\end{exercise}

\begin{exercise}
Implement the implicit midpoint rule for \eqref{eq:H} with Newton iteration and compare the energy error with RK4 for
$t\in[0,10^4]$ at fixed step size $h=10^{-2}$, once at low energy and once at high energy.
\end{exercise}

\begin{exercise}
Apply Kane's method with generalized speeds $u_r=\dot\theta_r$ ($r=1,2,3$) to a \emph{triple} pendulum
(masses $m_1,m_2,m_3$, rod lengths $l_1,l_2,l_3$). Show that the generalized mass matrix is
$M_{rs}=\sum_{k\ge\max(r,s)}m_k\,l_rl_s\cos(\theta_r-\theta_s)$ and obtain the three tensions with auxiliary speeds.
\end{exercise}

\begin{exercise}
Obtain the tension $T_1$ of the upper rod by introducing an auxiliary speed $u_4$ along $\bm e_1$ that displaces both
bobs, and show that the result agrees with the radial projection of \eqref{eq:newton1} on $\bm e_1$.
\end{exercise}

\begin{exercise}
Modify \texttt{midpoint\_step} in the Fortran code to use Newton's method with a finite-difference Jacobian, and
find the largest step for which each of the fixed-point and Newton versions converges for the chaotic initial
condition of Table~\ref{tab:energy}.
\end{exercise}

\begin{exercise}
Add a fourth-order Gauss--Legendre (two-stage) implicit Runge--Kutta integrator to the Fortran module and compare
its efficiency (energy error versus CPU time) with the Yoshida composition of Verlet steps.
\end{exercise}

\begin{thebibliography}{99}

\bibitem{goldstein} H.~Goldstein, C.~Poole and J.~Safko, \emph{Classical Mechanics}, 3rd ed.
(Addison-Wesley, 2002).

\bibitem{landau} L.~D. Landau and E.~M. Lifshitz, \emph{Mechanics}, 3rd ed. (Pergamon Press, Oxford, 1976).

\bibitem{arnold} V.~I. Arnold, \emph{Mathematical Methods of Classical Mechanics}, 2nd ed., Graduate Texts in
Mathematics 60 (Springer, 1989).

\bibitem{abrahammarsden} R.~Abraham and J.~E. Marsden, \emph{Foundations of Mechanics}, 2nd ed.
(Benjamin/Cummings, 1978).

\bibitem{kane1961} T.~R. Kane, ``Dynamics of nonholonomic systems,'' \emph{J. Appl. Mech.} \textbf{28} (1961) 574--578.

\bibitem{kanewang} T.~R. Kane and C.~F. Wang, ``On the derivation of equations of motion,''
\emph{J. Soc. Indust. Appl. Math.} \textbf{13} (1965) 487--492.

\bibitem{kanelevinson} T.~R. Kane and D.~A. Levinson, \emph{Dynamics: Theory and Applications}
(McGraw-Hill, 1985).

\bibitem{gibbs} J.~W. Gibbs, ``On the fundamental formulae of dynamics,'' \emph{Amer. J. Math.} \textbf{2} (1879) 49--64.

\bibitem{appell} P.~Appell, ``Sur une forme g\'en\'erale des \'equations de la dynamique,''
\emph{C. R. Acad. Sci. Paris} \textbf{129} (1899) 459--460.

\bibitem{baumgarte} J.~Baumgarte, ``Stabilization of constraints and integrals of motion in dynamical systems,''
\emph{Comput. Methods Appl. Mech. Eng.} \textbf{1} (1972) 1--16.

\bibitem{ziglin} S.~L. Ziglin, ``Branching of solutions and non-existence of first integrals in Hamiltonian mechanics. I,''
\emph{Funct. Anal. Appl.} \textbf{16} (1982) 181--189.

\bibitem{moralesruiz} J.~J. Morales-Ruiz, \emph{Differential Galois Theory and Non-Integrability of Hamiltonian
Systems} (Birkh\"auser, 1999).

\bibitem{mrr} J.~J. Morales-Ruiz and J.-P. Ramis, ``Galoisian obstructions to integrability of Hamiltonian systems,
I and II,'' \emph{Methods Appl. Anal.} \textbf{8} (2001) 33--95 and 97--111.

\bibitem{kolmogorov} A.~N. Kolmogorov, ``On the conservation of conditionally periodic motions under small perturbation
of the Hamiltonian,'' \emph{Dokl. Akad. Nauk SSSR} \textbf{98} (1954) 527--530.

\bibitem{arnoldkam} V.~I. Arnold, ``Proof of a theorem of A. N. Kolmogorov on the preservation of conditionally
periodic motions under a small perturbation of the Hamiltonian,'' \emph{Russian Math. Surveys} \textbf{18}(5) (1963) 9--36.

\bibitem{moser} J.~Moser, ``On invariant curves of area-preserving mappings of an annulus,''
\emph{Nachr. Akad. Wiss. G\"ottingen Math.-Phys. Kl. II} (1962) 1--20.

\bibitem{lichtenberg} A.~J. Lichtenberg and M.~A. Lieberman, \emph{Regular and Chaotic Dynamics}, 2nd ed.
(Springer, 1992).

\bibitem{benettin} G.~Benettin, L.~Galgani, A.~Giorgilli and J.-M. Strelcyn, ``Lyapunov characteristic exponents for
smooth dynamical systems and for Hamiltonian systems; a method for computing all of them. Part 1: Theory,''
\emph{Meccanica} \textbf{15} (1980) 9--20.

\bibitem{shinbrot} T.~Shinbrot, C.~Grebogi, J.~Wisdom and J.~A. Yorke, ``Chaos in a double pendulum,''
\emph{Am. J. Phys.} \textbf{60} (1992) 491--499.

\bibitem{levien} R.~B. Levien and S.~M. Tan, ``Double pendulum: An experiment in chaos,''
\emph{Am. J. Phys.} \textbf{61} (1993) 1038--1044.

\bibitem{stachowiak} T.~Stachowiak and T.~Okada, ``A numerical analysis of chaos in the double pendulum,''
\emph{Chaos, Solitons \& Fractals} \textbf{29} (2006) 417--422.

\bibitem{hnw} E.~Hairer, S.~P. N{\o}rsett and G.~Wanner, \emph{Solving Ordinary Differential Equations I: Nonstiff
Problems}, 2nd ed. (Springer, 1993).

\bibitem{ghw} E.~Hairer, C.~Lubich and G.~Wanner, \emph{Geometric Numerical Integration}, 2nd ed.
(Springer, 2006).

\bibitem{sanzserna} J.~M. Sanz-Serna and M.~P. Calvo, \emph{Numerical Hamiltonian Problems}
(Chapman \& Hall, 1994).

\bibitem{leimkuhler} B.~Leimkuhler and S.~Reich, \emph{Simulating Hamiltonian Dynamics} (Cambridge University Press, 2004).

\bibitem{marsdenwest} J.~E. Marsden and M.~West, ``Discrete mechanics and variational integrators,''
\emph{Acta Numerica} \textbf{10} (2001) 357--514.

\bibitem{yoshida} H.~Yoshida, ``Construction of higher order symplectic integrators,''
\emph{Phys. Lett. A} \textbf{150} (1990) 262--268.

\end{thebibliography}

\end{document}
